% ===============================================================
\chapter{Introduction}
% ===============================================================

\section{Natural languages}

\emph{Natural languages} are spoken by and used to communicate with people. They can be studied by being broken down into the following 4 components:

\begin{definition}\leavevmode
\begin{itemize}
  \item \emph{Syntax}: the study of how expressions combine to form grammatical units.
  \item \emph{Semantics}: the study of linguistic meaning.
  \item \emph{Pragmatics}: the study of how speakers use language.
  \item \emph{Phonology}: the study of how sounds combine to form units.
\end{itemize}
\end{definition}

\begin{note}
Noam Chomsky uses the phrase \emph{infinite use of finite means}, which expresses how finite language units can span an infinite combination of those units.
\end{note}

\section{The use/mention distinction}

\begin{itemize}
  \item If a word is \emph{used}, we are using the word in the sense of either meaning/pragmatics.
  \item If a word is \emph{mentioned}, we are referring to the word itself.
\end{itemize}

\section{Object language and meta language}

\begin{definition}\leavevmode
\begin{itemize}
  \item \emph{Object language}: the language being studied.
  \item \emph{Meta language}: the language used to study the object language.
\end{itemize}
\end{definition}

\section{Key definitions}

\begin{definition}\leavevmode
\begin{itemize}
  \item \emph{Synonymy}: two words are synonymous when they have the same meaning.
  \item \emph{Ambiguous}: one word has multiple meanings.
  \begin{itemize}
    \item \emph{Lexical ambiguity}: a sentence is ambiguous because of an ambiguous word.
    \item \emph{Structural ambiguity}: a sentence is ambiguous because a word in a sentence can be interpreted using multiple syntactic structures.
  \end{itemize}
  \item \emph{Antonymy}: two words have two meanings that are, in some respect, opposed to each other.
  \item \emph{Homonymy}: two words are written and pronounced in the same way but each have distinct meanings.
  \item \emph{Polysemy}: when two words are ambiguous and their meanings are extremely closely related.
\end{itemize}
\end{definition}

\section{Propositions}

\begin{remark}[This is an assumption]\leavevmode
\begin{itemize}
  \item A proposition is expressed/conveyed by a particular utterance.
  \item Propositions are not sentence types.
\end{itemize}
\end{remark}

\section{The type/token distinction}

\begin{definition}\leavevmode
\begin{itemize}
  \item \emph{Token}: a verbalised or written occurrence of a particular expression/sentence.
  \item \emph{Type}: expression/sentence types are literally different categorisations of tokens.
\end{itemize}
\end{definition}

\section{Truth-conditional semantics}

\begin{definition}
To know the meaning is to know the truth conditions of a sentence.
\end{definition}

\section{The Principle of Compositionality}

\begin{definition}[Principle of Compositionality]
The meaning of a sentence is given by the meaning of its component parts and how they are combined.
\end{definition}

Any theory of meaning that conforms to this requires the following (we use metavariables to describe the last two):
\begin{itemize}
  \item \emph{Vocabulary}: words of the language (words, literally).
  \item \emph{Syntax}: theory of how vocabulary forms to combine into grammatically acceptable sentences (describing what is a sentence). Closure means recursive language: sentences of the same language combine to form the same language.
  \item \emph{Semantics}: how the meanings of component parts combine (typically truth value, or describing semantic value).
\end{itemize}
